% 176_T0_Shor_Photon_De_ch.tex -- Kapitelinhalt Dok. 176
% Aufbereitung und Transformation: zwei getrennte Vorgänge
% Vollständige Neufassung August 2026 (R75)

\section{Aufbereitung und Transformation: zwei getrennte Vorgänge}

\subsection*{Zusammenfassung}

Shors Algorithmus zerfällt in zwei Vorgänge, die regelmäßig vermengt
werden, obwohl sie nichts miteinander zu tun haben:

\begin{center}
\begin{tabular}{lll}
\toprule
 & \textbf{Aufbereitung} & \textbf{Transformation} \\
\midrule
Aufgabe & Modulare Exponentiation & Periodenextraktion \\
Eingang & Zahlen $a$, $N$ & aufbereitetes Signal \\
Ausgang & Signal mit $f(x) = a^x \bmod N$ & Frequenzdarstellung \\
Natur & Arithmetisches Problem & Wellenoperation \\
Technologie & Klassischer Schaltkreis & Interferenz / Modulation \\
Aufwand & \textbf{Teuer} ($O(n^3)$ Gatter) & \textbf{Billig} ($O(n^2)$ Gatter) \\
\bottomrule
\end{tabular}
\end{center}

Die Aufbereitung ist schwierig, weil die modulare Exponentiation
schwierig ist -- in jeder Technologie gleichermaßen. Sie hängt
\emph{nicht} an quantenmechanischen Eigenschaften. Die Transformation
ist in jeder Technologie dieselbe Wellenoperation und in jeder billig.
Der Unterschied zwischen Optik und Qubit-Register liegt weder im einen
noch im anderen, sondern in Betriebstemperatur und
Kopplungsmechanismus.

\begin{keyresult}[Drei getrennte Parameterräume]
Zur Vermeidung von Verwechslungen unterscheidet dieses Dokument strikt:
\begin{enumerate}
  \item \textbf{Geometrischer Raum:} $\xi = 4/30000 \approx 1{,}333\times10^{-4}$
    -- der Strukturparameter der FFGFT. Fest, nicht einstellbar.
  \item \textbf{Algebraischer Raum:} $\xi_{\text{Higgs}} \approx
    1{,}038\times10^{-5}$ -- Feldkorrektur im Vakuumerwartungswert.
  \item \textbf{Zahlentheoretischer Raum} $\mathbb{Z}/N\mathbb{Z}$:
    \emph{kein} universeller Parameter. Die Periode $r$ wird durch
    $\mathrm{lcm}(p-1,q-1)$ bestimmt -- eine individuelle arithmetische
    Eigenschaft von $N$.
\end{enumerate}
Die Faktorisierung lebt ausschließlich im dritten Raum. Weder $\xi$ noch
$\xi_{\text{Higgs}}$ haben dort eine Funktion.
\end{keyresult}

\section{Die Aufbereitung}
\label{sec:aufbereitung}

\subsection{Was die Aufbereitung ist}

Die Aufbereitung ist die Berechnung der modularen Exponentiation
$f(x) = a^x \bmod N$. Sie ist ein \emph{reines arithmetisches Problem}.

Die Quelle~\cite{IACR2014} stellt klar, dass Shor die Aufbereitung in
Superposition \textbf{nicht spezifiziert} hat:
\begin{quote}
``P.~Shor has specified the operations for the process
$|0\rangle|0\rangle \rightarrow \frac{1}{\sqrt{q}}\sum_{a=0}^{q-1}|a\rangle|0\rangle$,
but \textbf{not specified the operations for the process}
$\frac{1}{\sqrt{q}}\sum_{a=0}^{q-1}|a\rangle|0\rangle \rightarrow
\frac{1}{\sqrt{q}}\sum_{a=0}^{q-1}|a\rangle|x^a\,(\bmod\,n)\rangle$.''
\cite{IACR2014}
\end{quote}
Shor spezifizierte \textbf{nur} den klassischen Prozess~\cite{Shor1997}:
\[
  (a, 1) \rightarrow (a,\; x^a \bmod N).
\]

Die Aufbereitung ist der \textbf{Flaschenhals} von Shors
Algorithmus~\cite{BSI2023, Padho2026}. Sie macht etwa
\textbf{95\,\% der gesamten Gatterkosten} aus, während die QFT nur
etwa 5\,\% ausmacht~\cite{Padho2026}.

\begin{keyresult}[Die Aufbereitung ist ein arithmetisches Problem]
Die Aufbereitung ist ein klassischer Prozess, spezifiziert als
$(a,1) \rightarrow (a, x^a \bmod N)$~\cite{Shor1997, IACR2014}.
Die modulare Exponentiation ist der \textbf{Flaschenhals} des
Algorithmus und macht rund 95\,\% der Gatterkosten
aus~\cite{BSI2023, Padho2026}.
\end{keyresult}

\subsection{Die Kosten der Aufbereitung}

Die modulare Exponentiation ist in \emph{jeder Technologie} gleich
aufwändig. Die Implementierung in Quantenschaltkreisen ist aufgrund
der notwendigen Reversibilität und Uncomputation sogar \textbf{teurer}
als die klassische Implementierung~\cite{TrappedIon2021, UoA2024}.

\begin{center}
\begin{tabular}{lll}
\toprule
Realisierung & Aufwand & für RSA-2048 \\
\midrule
Klassisches binäres System & $O(n^3)$ &
  Standard Square \& Multiply~\cite{Shor1997} \\
Quantenschaltkreis & $O(n^3)$ &
  Zusätzlicher Overhead durch Reversibilität, \\
  & & Uncomputation, Toffoli-Gatter~\cite{TrappedIon2021} \\
Hybrid (klass.\ Vorberechnung) & Reduziert &
  Um Faktor $\omega$ reduzierte Quantenmultiplikationen~\cite{Ekera2016} \\
Optische Maske & $2^n$ Elemente & $10^{617}$ Elemente \\
Klassische Tabelle & $2^n$ Werte & $10^{617}$ Werte \\
\bottomrule
\end{tabular}
\end{center}

\subsection{Was die Superposition dabei leistet}

\textbf{[X]} Nichts. Der Zustand
$\sum_x \ket{x}\ket{a^x \bmod N}$ enthält zwar alle Paare, aber eine
Messung liefert genau eines -- zufällig gezogen. Danach weiß man
so viel wie nach einer einzigen klassischen Auswertung. Die
Superposition vereinfacht die Arithmetik nicht; sie legt denselben
Schaltkreis über alle Eingaben.

\section{Die hybride Aufbereitung: Binäre Steuerung trifft analoge Welle}
\label{sec:hybrid}

Die reine binäre Implementierung der modularen Exponentiation erstickt in
einer Gatteranzahl von $O(n^3)$. Die rein analoge Implementierung scheitert
am thermischen Rauschen und am Modulo-Problem.

Ein hybrider Ansatz -- bestehend aus diskreter binärer Ansteuerung und
analoger Wellenmultiplikation -- nutzt die Stärken beider Welten:

\begin{itemize}
  \item \textbf{Das Binäre (die Ansteuerung):} Übernimmt die exakte,
    rauschfreie Diskretisierung der Modulo-Logik. Sie steuert, welche
    Phasenpfade freigeschaltet werden.
  \item \textbf{Das Analoge (die Welle):} Übernimmt die kontinuierliche,
    parallele Phasenakkumulation und Multiplikation über Interferenz und
    Mischerstrukturen in Echtzeit.
\end{itemize}

\subsection{Was physikalisch geschieht}

In einem solchen Hybridsystem führt die analoge Welle die Multiplikation
nicht über Spannungspegel aus, sondern über \emph{kohärente
Phasenüberlagerung}.

Wenn zwei analoge Lichtwellen oder HF-Signale mit den Phasen $\phi_1$
und $\phi_2$ zusammengemischt werden, entsteht durch die Wechselwirkung
ein Zustand, dessen Gesamtwelle beide Phaseninformationen zeitgleich
in sich trägt:
\[
  \Psi_{\text{gesamt}} \propto e^{i\phi_1} + e^{i\phi_2}
\]
oder bei gekoppelten Systemen:
\[
  \ket{\phi_1} \otimes \ket{\phi_2}.
\]

In der klassischen Elektrotechnik und Signalverarbeitung heißt dies
\emph{kohärente Summation}, \emph{Frequenzmischung} oder
\emph{Phasenmultiplexing}.

In der Sprache der Quantenphysik heißt dieselbe mathematische
Überlagerung von Phasenzuständen im Hilbertraum \emph{Superposition}
(und bei gekoppelten Subsystemen \emph{Verschränkung}).

\subsection{Warum die Physik hier neue Begriffe verwendet}

Die Quantenphysik nutzt den Begriff Superposition an dieser Stelle aus
zwei Gründen, die für die Hardware-Entwicklung entscheidend sind:

\begin{enumerate}
  \item \textbf{Die Diskretisierung durch einzelne Quanten (Photonen):}
    Sobald der analoge Signalpegel so weit abgesenkt wird, dass man mit
    einzelnen Photonen arbeitet, ist die analoge Welle nicht mehr eine
    kontinuierliche Spannung, sondern eine Wahrscheinlichkeitsverteilung.
    Die ``analoge Überlagerung'' wird zur Quanten-Superposition.
  \item \textbf{Der Ausweg aus dem Rausch-Problem:}
    Klassisch-analoge Mischer leiden unter thermischem Rauschen
    ($k \cdot T \cdot B$), das das Signal nach wenigen Kaskaden
    unlesbar macht. Ein photonisches/quantenmechanisches System nutzt
    die Diskretheit der Photonen, um trotz analoger Überlagerung über
    Quanten-Fehlerkorrektur ein rauschfreies Durchschalten zu ermöglichen.
\end{enumerate}

\subsection{Fazit zur hybriden Aufbereitung}

Die ``Superposition'' bei der Aufbereitung ist im Grunde der Versuch,
eine analoge, parallele Phasenmultiplikation durch ein binäres
Steuerkorsett so zu stabilisieren, dass sie nicht im analogen Rauschen
untergeht.

Dass die theoretische Physik diesen hybriden Zustand aus binärer
Steuerung und analoger Wellenüberlagerung als Superposition von
Rechenpfaden bezeichnet, ändert nichts an der schaltungstechnischen
Realität:

\begin{keyresult}[Hybride Aufbereitung = binär gesteuertes Interferenznetzwerk]
Die hybride Aufbereitung ist ein \emph{phased-array-artiges
Interferenz-Netzwerk}, das binär konditioniert wird. Die
``Superposition'' ist hier die physikalische Beschreibung der analogen
Wellenüberlagerung -- nicht ein quantenmechanischer Zauber.
\end{keyresult}

\subsection{Kosten der hybriden Aufbereitung}

\begin{center}
\begin{tabular}{lll}
\toprule
Implementierung & Gatteranzahl & Charakteristik \\
\midrule
Rein binär & $O(n^3)$ & Erstickt bei großen $n$ \\
Rein analog & --- & Scheitert am Rauschen \\
Hybrid (binär gesteuert, analoge Welle) & Reduziert & Phased-Array-Interferenz \\
\bottomrule
\end{tabular}
\end{center}

Die hybride Aufbereitung ist kein Quantenprozess. Sie ist ein
\emph{klassisches Interferenznetzwerk}, das durch binäre Steuerung
konditioniert wird -- und das thermische Rauschen durch den Betrieb
im Einzelphotonen-Regime umgeht.

\section{Die Fouriertransformation}
\label{sec:ft}

\subsection{Die Fouriertransformation ist eine Wellenoperation}

\textbf{[K]} Die Fouriertransformation ist eine \emph{Wellenoperation}.
Sie ist unabhängig von der Aufbereitung und in jeder Technologie
billig~\cite{Optica2021}. Sie interessiert sich nicht dafür, woher das
Signal stammt -- sie übersetzt eine Darstellung in eine andere. Neue
Information entsteht dabei nicht: Was im Eingangssignal fehlt,
erscheint auch in der Ausgabe nicht.

Die QFT ist der Schritt, der die Periodeninformation aus dem Zustand
extrahiert~\cite{Kumar2026}. Sie bewirkt, dass Amplituden konstruktiv
interferieren -- und zwar nur bei Werten, die Vielfache von $2^n/r$
sind~\cite{GitHub2026}.

\subsection{Die optische Fouriertransformation}

Die optische Fouriertransformation (OFT) bietet massive
Vorteile~\cite{Optica2021, Wang2022}:

\begin{itemize}
  \item \textbf{Geschwindigkeit:} Latenzen im Pikosekundenbereich,
    bestimmt durch die optische Laufzeit~\cite{Optica2021}.
  \item \textbf{Bandbreite:} TBit/s-Echtzeitverarbeitung mit
    vernachlässigbarem Energieverbrauch~\cite{Optica2021}.
  \item \textbf{Implementierung:} Dispersive Fourier-Transformation,
    integrierte photonische Chips oder Surface-Plasmon-Polariton-basierte
    FT~\cite{Optica2021}.
\end{itemize}

Die OFT ist eine \textbf{Wellenoperation}, realisiert durch Interferenz
und Phasenmodulation. Sie ist \textbf{kein rechenintensiver
Prozess}~\cite{Wang2022}.

\begin{keyresult}[Die OFT ist die schnelle Operation]
Die optische Fouriertransformation ist extrem schnell, energieeffizient
und benötigt keine Qubits~\cite{Optica2021, Wang2022}.
\end{keyresult}

\subsection{Die Kosten}

\begin{center}
\begin{tabular}{lll}
\toprule
Realisierung & Aufwand & Bemerkung \\
\midrule
Quanten-Fouriertransformation & $O(n^2)$ Gatter & sequenziell \\
Optisches MZI-Netz & Durchlaufzeit & simultan, Lichtgeschwindigkeit \\
Klassische FFT & $O(2^n \cdot n)$ & auf explizite Werte \\
\bottomrule
\end{tabular}
\end{center}

\subsection{Die entscheidende Erkenntnis}

Die Fouriertransformation -- ob optisch oder quantenhaft realisiert --
ist \emph{nicht das Problem}. Sie ist eine \textbf{schnelle, effiziente
Operation}, die durch Interferenz und Modulation realisiert wird.

Das \emph{eigentliche Problem} ist die \textbf{Aufbereitung} -- die
modulare Exponentiation. Sie ist ein arithmetisches Problem, das in
jeder Technologie gleich aufwändig ist.

\begin{keyresult}[Zwei völlig getrennte Vorgänge]
Die Aufbereitung und die Fouriertransformation sind \emph{zwei völlig
getrennte Vorgänge}. Sie haben nichts miteinander zu tun.

Die Fouriertransformation ist die \textbf{schnelle Operation}.
Die Aufbereitung ist das \textbf{echte Problem}.
\end{keyresult}

\subsection{Wozu sie dient}

Die gesuchte Periode $r$ ist eine \emph{globale} Eigenschaft der Folge:
Kein einzelner Funktionswert trägt sie; man erkennt sie erst im
Vergleich mehrerer Werte. Die Fouriertransformation macht diesen
Vergleich physikalisch -- durch Interferenz löscht sich alles aus
außer den Frequenzen $k/r$, und aus wenigen Messungen folgt $r$ per
Kettenbruchentwicklung.

\section{Photonische Realisierung}
\label{sec:photonik}

\subsection{Der thermische Vorteil optischer Realisierung}

\textbf{[K]} Ob ein Mode ohne Kühlung im Grundzustand ist, entscheidet das
Verhältnis $h\nu/(k_B T)$. Die mittlere thermische Besetzung beträgt
\begin{equation}
  \bar n = \frac{1}{e^{h\nu/k_B T} - 1}.
\end{equation}

\begin{center}
\begin{tabular}{lrrrr}
\toprule
System & $\nu$ (Hz) & $T$ (K) & $h\nu/k_BT$ & $\bar n$ \\
\midrule
Supraleitendes Qubit & $5{,}0\cdot10^{9}$ & $300$ & $8{,}0\cdot10^{-4}$
  & $1{,}25\cdot10^{3}$ \\
Supraleitendes Qubit & $5{,}0\cdot10^{9}$ & $0{,}020$ & $12{,}0$
  & $6{,}2\cdot10^{-6}$ \\
Photon 1550\,nm & $1{,}93\cdot10^{14}$ & $300$ & $30{,}9$
  & $3{,}7\cdot10^{-14}$ \\
Photon 600\,nm & $5{,}0\cdot10^{14}$ & $300$ & $80{,}0$
  & $1{,}8\cdot10^{-35}$ \\
\bottomrule
\end{tabular}
\end{center}

Bei $5$~GHz und Raumtemperatur sitzen rund \textbf{1250 thermische
Anregungen} im Mode -- das System ist vollständig durchmischt. Für eine
Besetzung unter $0{,}01$ ist $T < 52$~mK erforderlich. Bei optischen
Frequenzen ist der Grundzustand bei Raumtemperatur praktisch exakt
besetzt.

\textbf{Kühlkosten.} Die Carnot-Grenze verlangt mindestens
$3\cdot10^{4}$ Arbeitseinheiten je bei $10$~mK abgeführter
Wärmeeinheit; real deutlich mehr. Ein Verdünnungskryostat zieht
$10$--$25$~kW im Dauerbetrieb, ein optisches Netz nichts.

\begin{keyresult}[Thermodynamisch privilegiert, nicht rechnerisch]
Der thermische Vorteil folgt direkt aus $h\nu/k_BT$ und ist der
stärkste bekannte Vorteil photonischer gegenüber supraleitender
Realisierung. Er betrifft \emph{Infrastruktur} -- Kühlung, Energie,
Baugröße --, nicht die Rechenmächtigkeit.
\end{keyresult}

\subsection{Woher die exponentielle Dimension kommt}

\textbf{[K]} Entscheidend ist die Struktur des Zustandsraums, nicht das
Trägermedium:

\begin{center}
\begin{tabular}{lr}
\toprule
Aufbau & Dimension \\
\midrule
klassische Welle in $n$ Moden & $n$ \\
\emph{ein} Photon in $n$ Moden & $n$ \quad(dasselbe) \\
$n$ verschränkte Zweiniveausysteme & $2^n$ \\
$m$ ununterscheidbare Photonen in $n$ Moden & $\binom{n+m-1}{m}$ \\
\bottomrule
\end{tabular}
\end{center}

25 Photonen in 50 Moden ergeben $3{,}5\cdot10^{19}$. Der Sprung kommt
aus der \emph{Tensorproduktstruktur}, nicht aus dem Trägermedium.
Photonische Quantenrechner existieren (KLM-Schema, messungsbasierte
Verfahren).

\subsection{Was den Rechenvorteil trägt -- offene Frage}

\textbf{[K]} Der Satz von Gottesman-Knill (1998) zeigt: Jeder
Schaltkreis aus ausschließlich Clifford-Gattern ($H$, $S$, CNOT) ist
klassisch in Polynomzeit simulierbar, unabhängig von der erzeugten
Verschränkung. Verschränkung allein trägt den Rechenvorteil nicht.

Diskutiert werden Kontextualität, Nicht-Stabilisiertheit
(``Magie'') und die Interferenzstruktur selbst; keiner dieser
Kandidaten ist als alleiniger Träger etabliert. Dieses Dokument benennt
die Frage als offen.

\subsection{Topologische Einschränkung}

\textbf{[K]} Für geschlossene Trajektorien entscheidet die
\emph{logarithmische Windungszahl} $W = \log_2(r)$, nicht das
Frequenzverhältnis $r$:

\begin{center}
\begin{tabular}{llll}
\toprule
Fall & Verhältnis $r$ & Windungszahl $W$ & Bahn \\
\midrule
Reine Quinte & $3/2$ (rational) & $0{,}584963$ (irrational) & offen \\
Temperierte Quinte & $2^{7/12}$ (irrational) & $7/12$ (rational) & geschlossen \\
\bottomrule
\end{tabular}
\end{center}

\section{Strukturelle Grenzen}
\label{sec:weyl}

\subsection{Die Weyl-Obstruktion}

\textbf{[K]} Nach dem Weyl-Gesetz wächst das Laplace-Spektrum jeder
kompakten Mannigfaltigkeit nach $N(\lambda) \sim C_d\lambda^{d/2}$.
Arithmetische Zähldichten wachsen dagegen wie $N(T) \sim T\ln T$
(Riemann-von Mangoldt). Das ist keine Potenz.

\begin{keyresult}[Wichtige Abgrenzung]
Diese Obstruktion trifft \textbf{nicht} Shors Algorithmus. Shor sucht
\emph{eine} Periode \emph{einer konkreten} Zahl $N$ -- eine endliche
Aufgabe, für die ein Register mit $2n+3$ Qubits ausreicht.

Was Shor bei großen $N$ limitiert: Gatterzahl $\sim n^3$ und
Fehlerkorrektur-Overhead. Für RSA-2048: etwa $4099$ logische Qubits,
$8{,}6\times10^9$ Gatter, $4{,}1\times10^6$ physikalische Qubits.
Das ist eine \emph{Ressourcengrenze}, nicht eine strukturelle.
\end{keyresult}

\subsection{Primzahlen als Generatoren}

\textbf{[K]} Die tragende Resonanzstruktur beruht auf Primzahlen. Der
Bikohärenztest an 2469 Riemann-Nullstellen (Dok.~316) zeigt
Phasenkopplung genau bei Primzahlpotenzen ($b^2 = 0{,}76$--$0{,}85$)
und keine bei zusammengesetzten Zahlen ($b^2 = 0{,}018$--$0{,}025$).

\section{Einordnung des Quantenrechners}

\textbf{[K]} Shors Algorithmus faktorisiert auf Quantenhardware
theoretisch in $O((\log N)^3)$ -- eine reale Leistung für ein
spezifisches Problem.

\begin{itemize}
  \item \textbf{[X]} \emph{Allgemeiner exponentieller Quantenvorteil:}
    nicht nachgewiesen.
  \item \textbf{[X]} \emph{Vorteil liegt in der FT:} Diese ist eine
    Wellenoperation, in der Optik sogar schneller realisierbar. Der
    Flaschenhals ist die Aufbereitung.
  \item \textbf{[X]} \emph{Weyl limitiert Shor:} Shor löst eine
    endliche Aufgabe; Weyl gilt für unendliche Spektren.
\end{itemize}

\section{Bilanz}

\begin{center}
\begin{tabular}{ll}
\toprule
Aussage & Status \\
\midrule
Aufbereitung ist ein arithmetisches Problem & \textbf{[K]} \\
Aufbereitung ist der Flaschenhals (95\,\% der Gatterkosten) & \textbf{[K]} \\
Quanten-Aufbereitung ist teurer als klassische & \textbf{[K]} \\
Aufbereitung und FT sind unabhängig & \textbf{[K]} \\
Aufbereitung ist arithmetisch, nicht quantenmechanisch & \textbf{[K]} \\
FT ist eine Wellenoperation, in jeder Technologie billig & \textbf{[K]} \\
Optische FT ist extrem schnell (Pikosekunden, TBit/s) & \textbf{[K]} \\
Optik thermodynamisch privilegiert & \textbf{[K]} $h\nu/k_BT$ \\
Exponentielle Dimension aus dem Medium & \textbf{[X]} Tensorstruktur \\
Verschränkung trägt den Rechenvorteil & \textbf{[X]} Gottesman-Knill \\
Weyl-Obstruktion limitiert Shor & \textbf{[X]} Ressourcengrenze \\
Allgemeiner Quantenvorteil & \textbf{[X]} nicht nachgewiesen \\
$\xi$ hat Funktion in $\mathbb{Z}/N\mathbb{Z}$ & \textbf{[X]} getrennte Räume \\
\bottomrule
\end{tabular}
\end{center}

\section{Bezug zum Korpus}

Die Befunde stammen aus Dok.~316 (August 2026): Weyl-Obstruktion,
Bikohärenztest, Windungszahl gegen Frequenzverhältnis,
Fouriertransformation als Projektion. Dok.~161 und Dok.~162 behandeln
analoge Ising-Maschinen und physikalische Relaxation als Optimierungsansatz.
Die Überarbeitung ist unter R75 in Dok.~190 deklariert.

\section{Verifikation}

Die Aussagen dieses Dokuments sind durch drei Skripte reproduzierbar
(reine Standardbibliothek, Sollwerte als Assertions):

\begin{itemize}
  \item \texttt{s1\_periodensuche.py} -- Korrektheit der Periodenfindung,
    Skalierungsmessung, Primfaktorzerlegung der $\sigma$-Nenner,
    Verfahrensvergleich.
  \item \texttt{s2\_grenzen.py} -- Weyl-Obstruktion, Abgrenzung von
    Shors Aufgabe, Ressourcenabschätzung, Windungszahl-Probe.
  \item \texttt{s3\_thermik\_zustandsraum.py} -- thermische Besetzung nach
    Trägerfrequenz, Kühlkosten nach Carnot, Dimensionsvergleich,
    Gottesman-Knill.
\end{itemize}

\begin{thebibliography}{99}

\bibitem{Shor1997}
Shor, P.~W., ``Polynomial-Time Algorithms for Prime Factorization and
Discrete Logarithms on a Quantum Computer'',
\emph{SIAM Journal on Computing}, Vol.~26, No.~5, pp.~1484--1509 (1997).
DOI: 10.1137/S0097539795293172.
\url{https://arxiv.org/abs/quant-ph/9508027}

\bibitem{IACR2014}
``Remarks on Quantum Modular Exponentiation'',
\emph{IACR ePrint Archive}, 2014.
\url{https://eprint.iacr.org/archive/2014/828/}
Stellt klar, dass Shor die Aufbereitung in Superposition nicht
spezifizierte.

\bibitem{Padho2026}
``Modular Exponentiation as a Circuit'',
\emph{padho-wiki} (2026).
\url{https://learn.padho.ai/wiki/modular-exponentiation-circuit}
Zeigt: Aufbereitung macht 95\,\% der Gatterkosten aus, QFT nur 5\,\%.

\bibitem{BSI2023}
Bundesamt für Sicherheit in der Informationstechnik,
``Entwicklungsstand Quantencomputer'', Version~2.1 (2023).

\bibitem{Ekera2016}
Ekerå, M., ``Quantum information theory: Trading classical for quantum
computation using indirection'', Elsevier (2016).

\bibitem{TrappedIon2021}
``Time overhead of Shor's algorithm in trapped ion systems'',
\emph{International Journal of Theoretical Physics} (2021).

\bibitem{UoA2024}
``Quantum modular exponentiation circuit designs'',
\emph{University of Athens Thesis} (2024).

\bibitem{VanMeter2004}
Van Meter, R., ``Fast Quantum Modular Exponentiation'',
\emph{EQIS 2004 Poster}.

\bibitem{Optica2021}
``Optical Fourier transform for high-speed signal processing'',
\emph{Optica}, Vol.~8, No.~5 (2021).

\bibitem{Wang2022}
Wang, W., et al., ``Computing Shor's algorithmic steps with interference
patterns of classical light'',
\emph{Scientific Reports} 12, 21099 (2022).
DOI: 10.1038/s41598-022-25796-w.

\bibitem{Kumar2026}
Kumar, D., Raj, B., \& Singh, G.,
``Deconstructing Shor's Algorithm Using Quantum Fourier Transform'',
\emph{Indian Journal of Pure \& Applied Physics}, Vol.~64, No.~4 (2026).

\bibitem{GitHub2026}
``Shor's Algorithm -- Quantum Period Finding'',
\emph{GitHub Repository}, 2026.

\end{thebibliography}
